\documentclass[11pt]{article}
\usepackage{amsmath, amssymb, physics, mathtools, bm}
\usepackage{tikz, pgfplots}
\pgfplotsset{compat=1.18}
\usepackage[margin=2.3cm]{geometry}
\usepackage{siunitx, booktabs, hyperref}
\usepackage{braket}

\newcommand{\D}{\mathrm{D}}
\newcommand{\de}{\mathrm{d}}
\newcommand{\pd}[2]{\frac{\partial #1}{\partial #2}}
\newcommand{\pdd}[2]{\frac{\partial^{2} #1}{\partial #2^{2}}}
\newcommand{\Lap}{\nabla^{2}}

\title{B2 Symmetry and Relativity --- Model Solutions, 2020}
\author{}
\date{}

\begin{document}
\maketitle

Metric $(+,-,-,-)$ throughout. Capital bold $\bm U$, $\bm P$, etc.\ denote 4-vectors; lower-case bold $\vb v$, $\vb p$, etc.\ denote 3-vectors.

%==================================================================
\section*{Question 1 --- Threshold collisions and $\pi^0$ decay}

\textbf{Problem.} A particle of mass $m$, total energy $E$, hits a stationary mass $M$, producing $N$ particles of masses $m_i$. (a) Find threshold $E$. (b) Apply to $p+p\to p+p+\pi^0$: threshold proton kinetic energy and pion lab velocity. (c) Wavelength of the $\pi^0\to 2\gamma$ photons in the pion rest frame; show half the photons go into a lab cone of half-angle $\cos^{-1}(v/c)$. (d) How can the threshold be lowered.

\subsection*{(a)}

Total 4-momentum: $\bm P_\text{tot}=\bm P_m+\bm P_M$, conserved.

Lab frame, target at rest:
\[ \bm P_M=(Mc,\bm 0),\qquad \bm P_m=(E/c,\vb p). \]
Sum the 4-momenta:
\[ \bm P_\text{tot} = (E/c + Mc,\,\vb p). \]
Form the invariant square (metric $(+,-,-,-)$):
\[ \bm P_\text{tot}^2 = (E/c+Mc)^2-\vb p^{\,2}. \]
Expand the cross term $2(E/c)(Mc)=2EM$:
\[ \bm P_\text{tot}^2 = E^2/c^2 + 2EM + M^2c^2 - \vb p^{\,2}. \]
Apply on-shell projectile $E^2/c^2-\vb p^{\,2}=m^2c^2$:
\begin{equation}
\bm P_\text{tot}^2 = m^2c^2 + M^2c^2 + 2EM.
\label{eq:s-lab}
\end{equation}
At threshold all products are at rest in the CoM frame:
\[ \bm P_\text{tot}^2 = \Big(\textstyle\sum_i m_i\Big)^2 c^2. \]
Equate to \eqref{eq:s-lab}:
\[ \bigl(\textstyle\sum_i m_i\bigr)^2 c^2 = m^2c^2 + M^2c^2 + 2EM. \]
Isolate $2EM$:
\[ 2EM = \bigl[\bigl(\textstyle\sum_i m_i\bigr)^2 - m^2 - M^2\bigr]c^2. \]
Divide by $2M$:
\begin{equation}
\boxed{\;E_\text{min} = \frac{\bigl[\bigl(\sum_i m_i\bigr)^2 - m^2 - M^2\bigr]c^2}{2M}.\;}
\label{eq:Emin}
\end{equation}

\subsection*{(b)}

Set $m=M=m_p$, $\sum m_i = 2m_p+m_\pi$ in \eqref{eq:Emin}:
\[ E_\text{min} = \frac{\bigl[(2m_p+m_\pi)^2 - 2m_p^2\bigr]c^2}{2m_p}. \]
Expand $(2m_p+m_\pi)^2$:
\[ (2m_p+m_\pi)^2 = 4m_p^2 + 4m_p m_\pi + m_\pi^2. \]
Substitute and collect:
\[ E_\text{min} = \frac{\bigl[2m_p^2 + 4m_p m_\pi + m_\pi^2\bigr]c^2}{2m_p}. \]
Divide term by term:
\[ E_\text{min} = \Bigl(m_p + 2m_\pi + \frac{m_\pi^2}{2m_p}\Bigr)c^2. \]
Form kinetic energy $T = E - m_pc^2$:
\[ T_\text{min} = \Bigl(2m_\pi + \frac{m_\pi^2}{2m_p}\Bigr)c^2. \]
Plug in $m_\pi c^2 = \SI{135}{MeV}$, $m_p c^2 = \SI{938}{MeV}$:
\[ T_\text{min} = 2(135) + \frac{135^2}{2(938)}\ \text{MeV}. \]
Evaluate the second term:
\[ \frac{135^2}{2(938)} = \frac{18225}{1876} \approx 9.7\ \text{MeV}. \]
Sum the contributions:
\[ T_\text{min} = 270 + 9.7\ \text{MeV}. \]
\[ \boxed{\;T_\text{min} \approx \SI{280}{MeV} \approx 2.07\,m_\pi c^2.\;} \]
Comment: only $\sim 5\%$ above the naive $2m_\pi c^2$; the excess is the recoil kinetic energy of the products.

Pion lab velocity at threshold: all three products move together with the CoM velocity. Total 4-momentum:
\[ \bm P_\text{tot} = (E/c+m_pc,\,\vb p). \]
Magnitude of 3-momentum from on-shell projectile:
\[ |\vb p|c = \sqrt{E^2-m_p^2c^4}. \]
CoM velocity from $\vb v = c^2\vb p_\text{tot}/E_\text{tot}$:
\[ v_\pi = \frac{|\vb p|c^2}{E+m_pc^2}. \]
Total projectile energy $E = T_\text{min}+m_pc^2$:
\[ E = 280 + 938 = \SI{1218}{MeV}. \]
Compute $|\vb p|c$:
\[ |\vb p|c = \sqrt{1218^2 - 938^2}\ \text{MeV} = \SI{777}{MeV}. \]
Denominator:
\[ E+m_pc^2 = 1218 + 938 = \SI{2156}{MeV}. \]
Divide:
\[ \boxed{\;v_\pi/c = 777/2156 \approx 0.360.\;} \]

\subsection*{(c)}

In the pion rest frame, $\pi^0\to\gamma\gamma$ gives back-to-back photons. Energy conservation:
\[ E_\gamma' = m_\pi c^2/2. \]

Photon wavelength from $E_\gamma' = hc/\lambda'$:
\[ \lambda' = \frac{hc}{E_\gamma'} = \frac{2hc}{m_\pi c^2}. \]
Insert $hc = \SI{1240}{MeV.fm}$:
\[ \lambda' = \frac{2\cdot 1240}{135}\ \text{fm}. \]
Evaluate:
\[ \boxed{\;\lambda' \approx \SI{18.4}{fm}.\;} \]

Cone of half the photons. Isotropy in pion frame: half are emitted into the forward hemisphere $\theta'\in[0,\pi/2]$.

Aberration formula (boost along $+x$, $\theta'$ measured from $+x$ in pion frame):
\[ \cos\theta = \frac{\cos\theta'+\beta}{1+\beta\cos\theta'},\qquad \beta=v/c. \]
Set the hemisphere boundary $\theta'=\pi/2$:
\[ \cos\theta'=0. \]
Substitute:
\[ \cos\theta = \frac{0+\beta}{1+0} = \beta. \]
Invert:
\[ \boxed{\;\theta = \cos^{-1}(v/c).\;} \]
Half the photons (those with $\cos\theta'\ge 0$) lie inside this lab cone.

\subsection*{(d)}

Use a colliding-beam (symmetric) configuration: two protons of equal and opposite 3-momentum, so the lab \emph{is} the CoM frame.

Threshold in CoM frame: total 4-momentum has zero 3-momentum, so
\[ \bm P_\text{tot}^2 = (2E_\text{lab}/c)^2. \]
Set equal to the threshold invariant $(2m_p+m_\pi)^2 c^2$:
\[ (2E_\text{lab}/c)^2 = (2m_p+m_\pi)^2 c^2. \]
Solve for $E_\text{lab}$:
\[ E_\text{lab} = (m_p + m_\pi/2)c^2. \]
Kinetic energy per beam:
\[ T_\text{beam} = E_\text{lab} - m_p c^2 = \tfrac{1}{2}m_\pi c^2 = \SI{67.5}{MeV}. \]
Total kinetic energy (both beams):
\[ T_\text{total} = 2T_\text{beam} = m_\pi c^2 = \SI{135}{MeV}. \]
Compare to fixed-target value:
\[ \boxed{\;\Delta T = 280-135 = \SI{145}{MeV}\text{ saved; reduction factor }\approx 2.\;} \]
Final state: at threshold the three products $p,p,\pi^0$ are at rest in the lab.

%==================================================================
\section*{Question 2 --- Rocket dynamics and the photon rocket}

\textbf{Problem.} (a) Define $\bm U$, $\bm F$ and show $\de E/\de t = \vb v\cdot\vb f$. (b) Rocket with proper acceleration $a_0$: write equation of motion in rapidity $\eta$, derive trajectory $x^2-c^2t^2=(c^2/a_0)^2$, and show a photon emitted from origin never catches the rocket. (c) Photon rocket: show $m\to\alpha m$ with $\alpha=\sqrt{(1-\beta)/(1+\beta)}$; fuel fraction to reach $\gamma=8$ then decelerate; comment on massive-exhaust case.

\subsection*{(a)}

Definitions:
\[ \bm U = \frac{\de \bm X}{\de\tau} = \gamma(c,\vb v),\qquad \bm F = \frac{\de \bm P}{\de\tau} = \gamma(\dot E/c,\,\vb f), \]
where $\bm P=(E/c,\vb p)$, $\vb f=\de\vb p/\de t$ is the 3-force, $\tau$ proper time.

Pure force preserves rest mass: $\bm P^2 = m^2c^2 = $ const. Differentiate w.r.t.\ proper time $\tau$:
\[ 2\bm P\cdot\bm F = 0 \;\Rightarrow\; \bm P\cdot\bm F = 0. \]
Contract with the metric ($\bm P\cdot\bm F = P^0 F^0 - \vb p\cdot(\gamma\vb f)$):
\[ \frac{E}{c}\cdot\gamma\frac{\dot E}{c} - \vb p\cdot\gamma\vb f = 0. \]
Cancel the overall $\gamma$ and write $\dot E\equiv\de E/\de t$:
\[ \frac{E}{c^2}\frac{\de E}{\de t} - \vb p\cdot\vb f = 0. \]
Use $\vb p = (E/c^2)\vb v$:
\[ \frac{E}{c^2}\frac{\de E}{\de t} = \frac{E}{c^2}\,\vb v\cdot\vb f. \]
Cancel $E/c^2$:
\[ \boxed{\;\frac{\de E}{\de t} = \vb v\cdot\vb f.\;} \]

\subsection*{(b)}

Rapidity definition: $\cosh\eta=\gamma$, $\sinh\eta=\gamma\beta$, so $v/c=\tanh\eta$.

In the instantaneous rest frame the 4-acceleration is purely spatial with magnitude $a_0$:
\[ |A^\mu A_\mu| = a_0^2. \]
(Strictly spacelike with $(+,-,-,-)$, so $A^\mu A_\mu = -a_0^2$; we equate magnitudes.) Use the stated hint:
\[ c^2\dot\eta^2\cosh^2\eta = a_0^2. \]
Take the positive root:
\[ c\,\dot\eta\cosh\eta = a_0. \]
Recognise $\dot\eta\cosh\eta = \de(\sinh\eta)/\de t$:
\[ c\,\frac{\de}{\de t}\sinh\eta = a_0. \]
Integrate with $\eta(0)=0$:
\[ c\sinh\eta = a_0 t. \]
\begin{equation}
\sinh\eta = a_0 t/c.
\label{eq:eta-t}
\end{equation}
Lab velocity from $v = c\tanh\eta = c\sinh\eta/\cosh\eta$:
\[ v = \frac{a_0 t}{\sqrt{1+(a_0 t/c)^2}}. \]
Position from $\de x/\de t = v$:
\[ x(t) = \int_0^t \frac{a_0 t'\,\de t'}{\sqrt{1+(a_0 t'/c)^2}}. \]
Substitute $u = 1+(a_0 t'/c)^2$, $\de u = 2(a_0^2/c^2) t'\,\de t'$:
\[ x(t) = \frac{c^2}{2a_0}\int_1^{1+(a_0t/c)^2} u^{-1/2}\,\de u. \]
Evaluate:
\[ x(t) = \frac{c^2}{a_0}\bigl[\sqrt{1+(a_0 t/c)^2}-1\bigr]. \]
Choose origin at $x_0 = -c^2/a_0$ (shift $x\to x+c^2/a_0$):
\[ x(t) = \frac{c^2}{a_0}\sqrt{1+(a_0 t/c)^2}. \]
Square both sides:
\[ x^2 = (c^2/a_0)^2 + c^2 t^2. \]
Rearrange:
\begin{equation}
\boxed{\;x^2 - c^2 t^2 = (c^2/a_0)^2.\;}
\label{eq:hyperbola}
\end{equation}

\textbf{Photon never catches the rocket.} Photon worldline from $(t,x)=(0,0)$ in $+x$:
\[ x_\gamma = ct. \]
Rocket worldline from \eqref{eq:hyperbola}:
\[ x_\text{rocket}(t) = \sqrt{(c^2/a_0)^2 + c^2 t^2}. \]
Compare termwise:
\[ x_\text{rocket}^2 - x_\gamma^2 = (c^2/a_0)^2 > 0\quad\forall\,t, \]
so $x_\text{rocket}>x_\gamma$ for all $t\ge 0$. Asymptotic gap at large $t$ via $\sqrt{1+\epsilon}\approx 1+\epsilon/2$:
\[ x_\text{rocket} - ct = ct\Bigl[\sqrt{1+(c/a_0 t)^2}-1\Bigr] \approx \frac{c^3}{2 a_0^2\, t}\to 0^+. \]
The line $x=ct$ is the asymptote of the hyperbola: a \emph{Rindler horizon} bounding the rocket's past light cone of accessible signals from the origin.

\begin{center}
\begin{tikzpicture}[>=stealth, scale=1.0]
  \draw[->] (-0.3,0) -- (5.5,0) node[right] {$ct$};
  \draw[->] (0,-0.3) -- (0,4.5) node[above] {$x$};
  % light cone x = ct
  \draw[dashed] (0,0) -- (4.3,4.3) node[pos=0.85, sloped, above] {$x=ct$};
  % hyperbola x = sqrt(R^2 + (ct)^2), R=1.2
  \draw[thick, domain=0:4.2, samples=80, smooth, variable=\t]
       plot ({\t}, {sqrt(1.44 + \t*\t)});
  \node[circle,fill,inner sep=1.2pt,label=left:{$O$}] at (0,0) {};
  \node[circle,fill,inner sep=1.2pt,label=left:{$c^2/a_0$}] at (0,1.2) {};
  \node at (3.6,4.0) {rocket};
  \draw[->] (3.2,3.7) -- (2.7,3.45);
\end{tikzpicture}
\end{center}

\subsection*{(c)}

Initial state (lab frame): rocket of mass $m$ at rest, $\bm P_i = (mc,\bm 0)$.

After emission: rocket has mass $m'$, lab speed $v$ along $+\hat{\vb v}$, 4-momentum $\bm P_r=\gamma m'(c,\vb v)$. Photons travel backward:
\[ \bm P_\gamma = (E_\gamma/c,\,-(E_\gamma/c)\hat{\vb v}). \]

Energy conservation:
\[ mc = \gamma m'c + E_\gamma/c. \]
Momentum conservation along $\hat{\vb v}$:
\[ 0 = \gamma m' v - E_\gamma/c. \]
Solve the second for $E_\gamma/c$:
\[ E_\gamma/c = \gamma m' v. \]
Substitute into the first:
\[ mc = \gamma m'c + \gamma m' v. \]
Factor $\gamma m'$ and divide by $c$:
\[ m = \gamma m'(1+\beta). \]
Solve for $m'/m$:
\[ \frac{m'}{m} = \frac{1}{\gamma(1+\beta)}. \]
Use $1/\gamma = \sqrt{1-\beta^2} = \sqrt{(1-\beta)(1+\beta)}$:
\[ \frac{m'}{m} = \frac{\sqrt{(1-\beta)(1+\beta)}}{1+\beta}. \]
Simplify:
\begin{equation}
\boxed{\;\alpha \equiv \frac{m'}{m} = \sqrt{\frac{1-\beta}{1+\beta}}.\;}
\label{eq:alpha}
\end{equation}

\textbf{Fuel fraction.} Two stages of equal $|\beta_1|$: accelerate to $\beta_1$ (factor $\alpha_1$), then decelerate from $\beta_1$ to rest (factor $\alpha_1$ again, applied in the instantaneous rest frame of the rocket at the start of deceleration --- the derivation depends only on $|\beta|$, not its sign):
\[ m_\text{final}/m_\text{initial} = \alpha_1\cdot\alpha_1 = \alpha_1^2. \]
Compute $\beta_1$ from $\gamma_1=8$:
\[ \beta_1 = \sqrt{1-1/\gamma_1^2} = \sqrt{1-1/64}. \]
Evaluate:
\[ \beta_1 = \sqrt{63}/8 \approx 0.9922. \]
Form numerator and denominator:
\[ 1-\beta_1 \approx 0.00784,\qquad 1+\beta_1 \approx 1.9922. \]
Divide:
\[ \alpha_1^2 = \frac{1-\beta_1}{1+\beta_1} \approx \frac{0.00784}{1.9922} \approx 3.93\times 10^{-3}. \]
Fuel fraction:
\[ \boxed{\;f = 1 - \alpha_1^2 \approx 0.996,\text{ i.e.\ }\approx 99.6\%\text{ of initial mass is fuel}.\;} \]

\textbf{Massive exhaust.} Photons carry maximal momentum per unit energy ($p=E/c$). Massive exhaust at the same energy carries $p<E/c$, so a larger expended rest-mass is needed for the same impulse: the fuel fraction increases.

%==================================================================
\section*{Question 3 --- Field tensor, gauge invariance, moving wire}

\textbf{Problem.} (a) Components of $F^{\mu\nu}$; show $\partial_\nu F^{\mu\nu}=\mu_0 J^\mu$ encodes Gauss + Amp\`ere; build an invariant of $\vb E,\vb B$. (b) From $F^{\mu\nu}=\partial^\mu A^\nu-\partial^\nu A^\mu$ show gauge invariance and reduce Maxwell to $\partial_\nu\partial^\nu A^\mu=-\mu_0 J^\mu$ in a suitable gauge. (c) Charged wire moving along its length: find $\vb E,\vb B$ in rest frame $S$ and lab $S'$; is there a frame with $\vb E=0$?

\subsection*{(a)}

Field tensor (metric $(+,-,-,-)$, raised indices; standard convention):
\[ F^{\mu\nu} = \begin{pmatrix} 0 & -E_x/c & -E_y/c & -E_z/c \\ E_x/c & 0 & -B_z & B_y \\ E_y/c & B_z & 0 & -B_x \\ E_z/c & -B_y & B_x & 0 \end{pmatrix}, \]
i.e.\ $F^{0i} = -E_i/c$, $F^{ij} = -\epsilon^{ijk}B_k$. 4-current: $J^\mu=(c\rho,\vb J)$. (The question's equation $\partial_\nu F^{\mu\nu}=\mu_0 J^\mu$ contracts on the second index; by antisymmetry $\partial_\nu F^{\mu\nu} = -\partial_\nu F^{\nu\mu}$, equivalent up to a sign to the canonical Steane form $\partial_\mu F^{\mu\nu}=\mu_0 J^\nu$. Both yield Maxwell --- we work with the canonical form below.)

Take $\nu=0$ in $\partial_\mu F^{\mu\nu}=\mu_0 J^\nu$:
\[ \partial_\mu F^{\mu 0} = \mu_0 J^0. \]
Antisymmetry kills $\mu=0$:
\[ \partial_i F^{i0} = \mu_0 c\rho. \]
Insert $F^{i0}=E_i/c$:
\[ \frac{1}{c}\nabla\cdot\vb E = \mu_0 c\rho. \]
Use $\mu_0 c^2 = 1/\epsilon_0$:
\[ \nabla\cdot\vb E = \rho/\epsilon_0. \quad\checkmark\text{(Gauss)} \]

Take $\nu=i$ (spatial):
\[ \partial_\mu F^{\mu i} = \partial_0 F^{0i}+\partial_j F^{ji} = \mu_0 J^i. \]
Insert $F^{0i}=-E_i/c$ and $F^{ji}=-F^{ij}=+\epsilon^{ijk}B_k$:
\[ -\frac{1}{c^2}\partial_t E_i + \epsilon^{ijk}\partial_j B_k = \mu_0 J^i. \]
Identify $\epsilon^{ijk}\partial_j B_k = (\nabla\times\vb B)_i$:
\[ (\nabla\times\vb B)_i - \frac{1}{c^2}\partial_t E_i = \mu_0 J^i. \]
Promote to vector form:
\[ \nabla\times\vb B - \frac{1}{c^2}\pd{\vb E}{t} = \mu_0\vb J. \quad\checkmark\text{(Amp\`ere--Maxwell)} \]

\textbf{Invariant.} Form the scalar $F_{\mu\nu}F^{\mu\nu}$ by lowering indices with $\eta = \mathrm{diag}(+,-,-,-)$:
\[ F_{\mu\nu}F^{\mu\nu} = 2\sum_i F^{0i}F_{0i} + \sum_{i,j} F^{ij}F_{ij}. \]
Compute the time--space piece (one index lowered flips sign): with $F^{0i}=E_i/c$, $F_{0i}=-E_i/c$,
\[ \sum_i F^{0i}F_{0i} = -E^2/c^2. \]
Compute the space--space piece using $F^{ij}=-\epsilon^{ijk}B_k$ and $F_{ij}=F^{ij}$:
\[ \sum_{i,j} F^{ij}F_{ij} = 2B^2. \]
Sum the contributions:
\[ F_{\mu\nu}F^{\mu\nu} = -2E^2/c^2 + 2B^2. \]
Halve to isolate the invariant:
\[ \boxed{\;\tfrac{1}{2}F_{\mu\nu}F^{\mu\nu} = B^2 - E^2/c^2 = \text{Lorentz invariant}.\;} \]

\subsection*{(b)}

Gauge transformation $A^\mu\to A^\mu+\partial^\mu\chi$ for any scalar $\chi$:
\[ F^{\mu\nu}\to \partial^\mu(A^\nu+\partial^\nu\chi)-\partial^\nu(A^\mu+\partial^\mu\chi). \]
Distribute the derivatives:
\[ F^{\mu\nu}\to F^{\mu\nu} + \partial^\mu\partial^\nu\chi - \partial^\nu\partial^\mu\chi. \]
The mixed partials commute:
\[ \partial^\mu\partial^\nu\chi - \partial^\nu\partial^\mu\chi = 0. \]
Therefore:
\[ F^{\mu\nu}\to F^{\mu\nu}. \quad\checkmark \]
Equivalent 3-vector check: with $\phi\to\phi+\partial_t\chi$ and $\vb A\to\vb A - \nabla\chi$,
\[ \vb B = \nabla\times\vb A \to \nabla\times\vb A - \nabla\times\nabla\chi = \vb B. \]
\[ \vb E = -\nabla\phi-\partial_t\vb A \to -\nabla\phi - \partial_t\nabla\chi - \partial_t\vb A + \partial_t\nabla\chi = \vb E. \]

Substitute $F^{\mu\nu}=\partial^\mu A^\nu-\partial^\nu A^\mu$ into $\partial_\nu F^{\mu\nu}=\mu_0 J^\mu$:
\[ \partial_\nu\partial^\mu A^\nu - \partial_\nu\partial^\nu A^\mu = \mu_0 J^\mu. \]
Swap derivative order in the first term (partials commute):
\[ \partial^\mu(\partial_\nu A^\nu) - \partial_\nu\partial^\nu A^\mu = \mu_0 J^\mu. \]
Impose Lorenz gauge:
\[ \partial_\nu A^\nu = 0. \]
The first term vanishes, leaving:
\[ \boxed{\;\partial_\nu\partial^\nu A^\mu = -\mu_0 J^\mu.\;} \]
Note: $\partial_\nu\partial^\nu = (1/c^2)\partial_t^2 - \nabla^2 = \Box$, so this is $\Box A^\mu = -\mu_0 J^\mu$; the more common form $\Box A^\mu = \mu_0 J^\mu$ corresponds to the alternate sign convention on $F^{\mu\nu}$ noted in (a).

\subsection*{(c)}

Wire rest frame $S$: linear charge density $\rho$, no current. Apply Gauss to a cylinder of radius $r$ and length $L$:
\[ 2\pi r L\,E_r = \rho L/\epsilon_0. \]
Solve for $E_r$:
\[ E_r = \frac{\rho}{2\pi\epsilon_0 r}. \]
\[ \boxed{\;\vb E_S = \frac{\rho}{2\pi\epsilon_0 r}\hat{\vb r},\qquad \vb B_S = \bm 0.\;} \]

Lab frame $S'$: the wire moves at $\vb v = v\hat{\vb z}$ in $S'$, so $S'$ moves at $-\vb v = -v\hat{\vb z}$ relative to $S$. Field transformation for $S'$ moving with velocity $\vb u$ relative to $S$ ($\vb u=-\vb v$):
\[ \vb E'_\perp = \gamma(\vb E_\perp + \vb u\times\vb B_\perp),\qquad \vb B'_\perp = \gamma\bigl(\vb B_\perp - (\vb u/c^2)\times\vb E_\perp\bigr). \]
Insert $\vb B = \bm 0$ and $\vb u = -\vb v$:
\[ \vb E' = \gamma\vb E = \frac{\gamma\rho}{2\pi\epsilon_0 r}\hat{\vb r}. \]
\[ \vb B' = -\gamma(\vb u/c^2)\times\vb E = +\gamma(\vb v/c^2)\times E_r\hat{\vb r}. \]
Use $\vb v = v\hat{\vb z}$ and $\hat{\vb z}\times\hat{\vb r}=\hat{\bm\varphi}$:
\[ \vb B' = \frac{\gamma\rho v}{2\pi\epsilon_0 c^2 r}\hat{\bm\varphi}. \]
Apply $1/(\epsilon_0 c^2) = \mu_0$:
\[ \vb B' = \frac{\mu_0\gamma\rho v}{2\pi r}\hat{\bm\varphi}. \]
Identify lab-frame line density $\rho'=\gamma\rho$ (length contraction) and current $I'=\rho' v=\gamma\rho v$:
\[ \vb E' = \frac{\rho'}{2\pi\epsilon_0 r}\hat{\vb r},\qquad \vb B' = \frac{\mu_0 I'}{2\pi r}\hat{\bm\varphi}. \]
\[ \boxed{\;\vb E' \text{ Coulomb of } \rho'=\gamma\rho;\quad \vb B' \text{ Amp\`ere of } I' = \gamma\rho v.\;} \]

\textbf{Frame with $\vb E=0$?} Use the invariant $I_1 \equiv B^2 - E^2/c^2$. Evaluate in $S$:
\[ I_1\big|_S = 0 - E_S^2/c^2 = -E_S^2/c^2 < 0. \]
Lorentz invariance forces $I_1<0$ in every frame:
\[ B^2 - E^2/c^2 < 0 \;\Rightarrow\; E^2 > c^2 B^2 \ge 0\;\Rightarrow\; \vb E\ne\bm 0. \]
\[ \boxed{\;\text{No frame has }\vb E=0.\;} \]

%==================================================================
\section*{Question 4 --- EM wave 4-potential, intervals, photon CoM, relative speed}

\textbf{Problem.} (a) Write $A^\mu$ for a vacuum EM wave; show phase $\phi$ is invariant; show an isolated electron cannot emit a photon. (b) Condition for spacelike separation; can simultaneity be achieved? (c) Two photons with 4-momenta $\bm P_1$ ($+x$), $\bm P_2$ ($+y$), same $\omega$: rest mass and CoM velocity. (d) Derive the relative speed formula.

\subsection*{(a)}

4-potential of a vacuum plane wave:
\[ A^\mu(\bm X) = A_0^\mu \cos(\bm K\cdot\bm X), \]
with constant polarisation $A_0^\mu$, 4-wavevector $\bm K=(\omega/c,\vb k)$ (null: $\bm K^2=0$ in vacuum) and event $\bm X=(ct,\vb x)$.

Write the phase using the metric:
\[ \phi = \bm K\cdot\bm X = K_\mu X^\mu = \omega t - \vb k\cdot\vb x. \]
$K_\mu X^\mu$ is a contraction of two 4-vectors:
\[ \boxed{\;\phi = K_\mu X^\mu \text{ is Lorentz invariant}.\;} \]

\textbf{Electron $\to e + \gamma$ forbidden.} 4-momentum conservation:
\[ \bm P_e = \bm P_e' + \bm P_\gamma. \]
Isolate the photon and square:
\[ \bm P_\gamma = \bm P_e - \bm P_e'. \]
Take the invariant square:
\[ \bm P_\gamma^2 = \bm P_e^2 - 2\bm P_e\cdot\bm P_e' + \bm P_e'^2. \]
Insert $\bm P_e^2 = \bm P_e'^2 = m_e^2c^2$ and $\bm P_\gamma^2 = 0$:
\[ 0 = 2m_e^2c^2 - 2\bm P_e\cdot\bm P_e'. \]
Solve for the dot product:
\[ \bm P_e\cdot\bm P_e' = m_e^2 c^2. \]
Evaluate in the initial electron rest frame, $\bm P_e = (m_ec,\bm 0)$, $\bm P_e' = (E_e'/c,\vb p_e')$:
\[ \bm P_e\cdot\bm P_e' = m_e c\cdot E_e'/c = m_e E_e'. \]
Equate:
\[ m_e E_e' = m_e^2 c^2. \]
Cancel one $m_e$:
\[ E_e' = m_e c^2. \]
But $E_e' = \sqrt{m_e^2 c^4 + |\vb p_e'|^2 c^2}\ge m_e c^2$, with equality only at $\vb p_e'=\bm 0$. Then $\bm P_\gamma = \bm P_e-\bm P_e' = 0$, i.e.\ no photon:
\[ \boxed{\;\text{No photon can be emitted; }e^-\to e^-\gamma\text{ forbidden in vacuum}.\;} \]

\subsection*{(b)}

Squared interval with $(+,-,-,-)$:
\[ (\Delta\bm X)^2 = c^2(\Delta t)^2 - (\Delta\vb x)^2. \]
Spacelike means $(\Delta\bm X)^2 < 0$:
\[ \boxed{\;c^2(t_2-t_1)^2 < (\vb x_2-\vb x_1)^2.\;} \]

Simultaneity in $S'$: require $\Delta t' = 0$. Lorentz boost along $\hat{\vb n}$ with speed $u$:
\[ c\Delta t' = \gamma\bigl(c\Delta t - (u/c)\,\hat{\vb n}\cdot\Delta\vb x\bigr). \]
Set $\Delta t' = 0$:
\[ c\Delta t = (u/c)\,\hat{\vb n}\cdot\Delta\vb x. \]
Choose $\hat{\vb n}$ along $\Delta\vb x$:
\[ c\Delta t = (u/c)|\Delta\vb x|. \]
Solve for $u$:
\[ u = \frac{c^2\,\Delta t}{|\Delta\vb x|}. \]
Spacelike $\Rightarrow c|\Delta t| < |\Delta\vb x| \Rightarrow |u| < c$, achievable:
\[ \boxed{\;\text{Yes, a frame }S'\text{ exists in which the events are simultaneous}.\;} \]

\subsection*{(c)}

Photon 4-momenta (using $E=\hbar\omega$, $|\vb p|c = E$):
\[ \bm P_1 = (\hbar\omega/c)(1,1,0,0),\qquad \bm P_2 = (\hbar\omega/c)(1,0,1,0). \]
Sum component-wise:
\[ \bm P = \bm P_1+\bm P_2 = (\hbar\omega/c)(2,1,1,0). \]
Square with $(+,-,-,-)$:
\[ \bm P^2 = (\hbar\omega/c)^2\bigl(2^2-1^2-1^2-0^2\bigr). \]
Simplify:
\[ \bm P^2 = 2(\hbar\omega/c)^2. \]
Identify rest mass via $\bm P^2 = M^2 c^2$:
\[ M^2 c^2 = 2(\hbar\omega/c)^2. \]
Solve:
\[ \boxed{\;M = \sqrt{2}\,\hbar\omega/c^2.\;} \]

CoM velocity from $\vb v_\text{CoM} = c^2\vb p_\text{tot}/E_\text{tot}$:
\[ \vb p_\text{tot} = (\hbar\omega/c)(\hat{\vb x}+\hat{\vb y}),\qquad E_\text{tot} = 2\hbar\omega. \]
Substitute:
\[ \vb v_\text{CoM} = \frac{c^2 (\hbar\omega/c)(\hat{\vb x}+\hat{\vb y})}{2\hbar\omega}. \]
Cancel $\hbar\omega$:
\[ \vb v_\text{CoM} = \frac{c}{2}(\hat{\vb x}+\hat{\vb y}). \]
Magnitude:
\[ \boxed{\;\vb v_\text{CoM} = \tfrac{c}{2}(\hat{\vb x}+\hat{\vb y}),\quad |\vb v_\text{CoM}|=c/\sqrt{2}.\;} \]

\subsection*{(d)}

Particle 1 has 3-velocity $\vb u$ in $S$; particle 2 has velocity $\vb v$ in $S$. Boost from $S$ to particle 2's rest frame $S'$ (which moves at $\vb v$ in $S$); the relative speed is $w=|\vb u'|$.

Decompose $\vb u$ into components parallel and perpendicular to $\vb v$. Velocity transformation:
\[ \vb u'_\parallel = \frac{\vb u_\parallel - \vb v}{1-\vb u\cdot\vb v/c^2},\qquad \vb u'_\perp = \frac{\vb u_\perp}{\gamma_v(1-\vb u\cdot\vb v/c^2)}. \]
Abbreviate $D = 1-\vb u\cdot\vb v/c^2$.

Square magnitude (parallel and perpendicular pieces independent):
\[ |\vb u'|^2 = |\vb u'_\parallel|^2 + |\vb u'_\perp|^2. \]
Insert the transformations:
\[ |\vb u'|^2 = \frac{(\vb u_\parallel-\vb v)^2}{D^2} + \frac{\vb u_\perp^2}{\gamma_v^2 D^2}. \]
Factor out $1/D^2$:
\[ |\vb u'|^2 = \frac{1}{D^2}\Bigl[(\vb u_\parallel-\vb v)^2 + \vb u_\perp^2/\gamma_v^2\Bigr]. \]
Use $1/\gamma_v^2 = 1-v^2/c^2$:
\[ |\vb u'|^2 = \frac{(\vb u_\parallel-\vb v)^2 + (1-v^2/c^2)\vb u_\perp^2}{D^2}. \]
Expand the perpendicular piece:
\[ (\vb u_\parallel-\vb v)^2 + (1-v^2/c^2)\vb u_\perp^2 = (\vb u_\parallel-\vb v)^2 + \vb u_\perp^2 - (v^2/c^2)\vb u_\perp^2. \]
Combine $(\vb u_\parallel-\vb v)^2 + \vb u_\perp^2 = (\vb u-\vb v)^2$ (since $\vb v$ is along $\hat\parallel$):
\[ (\vb u_\parallel-\vb v)^2 + (1-v^2/c^2)\vb u_\perp^2 = (\vb u-\vb v)^2 - (v^2/c^2)\vb u_\perp^2. \]
Apply the hint $(\vb a\times\vb b)^2 = a^2 b^2 - (\vb a\cdot\vb b)^2$ with $\vb a=\vb v$, $\vb b=\vb u$:
\[ (\vb v\times\vb u)^2 = v^2 u^2 - (\vb u\cdot\vb v)^2. \]
Note $\vb u\cdot\vb v = u_\parallel v$ and $\vb u_\perp^2 = u^2 - u_\parallel^2$:
\[ v^2\vb u_\perp^2 = v^2 u^2 - v^2 u_\parallel^2 = v^2 u^2 - (\vb u\cdot\vb v)^2 = (\vb u\times\vb v)^2. \]
Substitute back:
\[ |\vb u'|^2 = \frac{(\vb u-\vb v)^2 - (\vb u\times\vb v)^2/c^2}{D^2}. \]
Take the positive square root:
\[ \boxed{\;w = |\vb u'| = \frac{\sqrt{(\vb u-\vb v)^2 - (\vb u\times\vb v)^2/c^2}}{1-\vb u\cdot\vb v/c^2}.\;} \]
Sanity: parallel motion $\vb u\times\vb v=0$ recovers $w=|u-v|/(1-uv/c^2)$; non-relativistic limit $w\to|\vb u-\vb v|$.

\end{document}
